\chapter{总结与展望}

\section{全文总结}

本文围绕面向仿真-实测迁移的跨域SAR目标识别问题展开研究，针对实测样本稀缺、域偏移显著以及工程部署受限等关键挑战，按照“面向仿真与少量实测数据训练”和“面向纯仿真数据训练”两类典型场景，分别研究了域适应与域泛化方法，并对所提出方法开展了轻量化部署验证。全文的主要研究工作与结论总结如下：

（1）针对面向仿真与少量实测数据训练场景下实测样本稀缺、跨域分布难以有效对齐的问题，第二章提出了一种基于像素级域对齐和类混淆损失的域适应SAR目标识别方法。该方法从像素级和特征级两个层面协同缓解域偏移：在像素级引入基于对比学习的非配对图像翻译模型CUT，使仿真图像在少样本约束下实现域间外观迁移并保持目标结构与散射属性；在特征级设计类混淆损失与一致性正则化机制，结合概率校准和不确定性加权，直接优化实测域判别边界。SAMPLE数据集实验结果表明，所提方法在目标域全样本（all）设置下达到99.15\%准确率，在极具挑战的1-shot条件下仍取得92.45\%准确率，体现出从数据充足到极少样本场景下较稳定的跨域识别能力，说明该方法能够在降低实测数据依赖的同时有效提升面向仿真-实测迁移的跨域识别性能。

（2）针对面向纯仿真数据训练场景下训练阶段实测域完全不可见、模型易过拟合仿真域的问题，第三章提出了一种基于多层级数据增强和对抗学习的域泛化SAR目标识别方法。该方法构建了由像素级数据增强、特征级风格混淆和域对抗学习组成的协同框架，从输入分布扩展、深层特征风格扰动以及域不变特征约束三个层面抑制域偏移。其中，像素级数据增强结合SAR图像特点扩展仿真数据分布，特征级风格混淆模拟多样化域风格，域对抗学习强化模型对域不变信息的提取能力。SAMPLE和S2M-5数据集实验结果表明，所提方法分别取得了96.93\%和93.70\%的准确率，跨数据集性能降幅仅为3.23\%，说明该方法在未知实测域上具有较好的识别精度和跨域稳定性，验证了其在纯仿真训练条件下解决仿真-实测迁移问题的有效性。

（3）为验证所提跨域SAR目标识别方法的工程应用可行性，第四章开展了面向国产AI芯片平台的轻量化部署研究。该章以前两章方法为基础，构建了融合知识蒸馏、后训练量化与板端推理验证的轻量化部署流程，通过教师—学生网络知识迁移和INT8量化压缩，在较好保持跨域识别性能的同时显著降低了模型规模与推理时延，并在华为Atlas 200I DK A2平台完成部署验证。实验结果表明，轻量化后的MobileNetV2模型体积压缩至6.67MB，C++单图推理时延降低至1.33ms；在第二章和第三章对应部署分支中，模型准确率下降幅度总体控制在约2-3\%。上述工作将本文前两章的方法创新进一步扩展为可压缩、可迁移、可部署的工程化成果，为跨域SAR目标识别模型的国产化应用提供了实现路径。

综上所述，本文面向仿真-实测迁移中的跨域SAR目标识别需求，围绕两类典型训练场景分别建立了域适应和域泛化方法，并通过轻量化部署验证了相关方法的工程应用可行性。相关结果表明，第二章方法在all和1-shot设置下分别达到99.15\%和92.45\%准确率，第三章方法在SAMPLE和S2M-5数据集上分别达到96.93\%和93.70\%准确率，第四章进一步在将模型体积压缩至6.67MB、单图推理时延降低至1.33ms的条件下较好保持了跨域识别性能。全文工作形成了一套面向仿真-实测迁移的跨域SAR目标识别方法体系，为实测样本受限条件下SAR目标识别模型的设计、训练与应用提供了理论方法和技术支撑。


\section{未来展望}

尽管本文围绕面向仿真-实测迁移的跨域SAR目标识别开展了较为系统的研究，并在方法设计与部署验证方面取得了一定成果，但受限于研究时间和实验条件，仍存在若干值得进一步深入探索的方向：

（1）\textbf{面向开放场景的跨域SAR目标识别方法研究。}
本文的域适应和域泛化方法均假设源域与目标域共享相同的类别空间。
然而在实际应用中，目标域可能出现训练阶段未涵盖的新目标类别（开集识别），
或源域与目标域的类别仅部分重叠（部分域适应）。
未来可结合开放集识别、广义零样本学习等技术，
研究类别空间不一致条件下的跨域SAR目标识别方法，
进一步提升模型在开放世界场景下的实用性。

（2）\textbf{基于大规模预训练模型的跨域SAR目标识别方法研究。}
近年来，视觉基础大模型（如ViT、SAM等）在自然图像领域展现了强大的跨域泛化能力。
未来可探索将大规模预训练视觉模型引入SAR目标识别领域，
通过参数高效微调（如LoRA、Adapter等）或提示学习（Prompt Learning）策略，
利用预训练模型中蕴含的丰富视觉先验知识增强SAR目标识别的跨域泛化性能，
同时降低对SAR领域大规模标注数据的依赖。

（3）\textbf{多模态融合的跨域SAR目标识别方法研究。}
当前本文仅利用单一的SAR幅度图像进行目标识别。
实际SAR系统可获取丰富的多模态信息，包括极化信息、相位信息以及电磁散射特性等。
未来可研究融合SAR多极化、多频段信息以及电磁散射先验知识的跨域识别方法，
充分利用SAR数据固有的物理特性以构建更具鲁棒性的域不变特征表示，
提升模型在复杂域偏移条件下的识别性能。

（4）\textbf{面向实际部署场景的端云协同推理架构研究。}
本文的部署验证主要面向单一边缘设备的独立推理场景。
在实际战场应用中，SAR目标识别系统往往需要在多级计算节点
（如机载终端、地面站、云端服务器）之间进行协同工作。
未来可研究面向跨域SAR目标识别的端云协同推理架构，
根据任务复杂度、通信带宽和时延约束动态分配计算负载，
在保证识别实时性的同时充分利用云端算力提升复杂场景下的识别精度，
推动相关技术在实际作战体系中的进一步应用。
